% Source-derived chapter 01: Introduction
% Original source: bounds-stalks-perverse-sheaves-characteristic-p-shende
\chapter{Introduction}
\label{bounds-stalks-perverse-sheaves-characteristic-p-shende-chapter-01}
\source{bounds-stalks-perverse-sheaves-characteristic-p-shende}

\begin{remark}[Source section]
\label{bounds-stalks-perverse-sheaves-characteristic-p-shende-chapter-01-source}
\source{bounds-stalks-perverse-sheaves-characteristic-p-shende}
\source{bounds-stalks-perverse-sheaves-characteristic-p-shende}
This chapter reproduces the corresponding section of the retrieved
LaTeX source, including its definitions, statements, examples, and
proofs. It has not yet been translated into Lean.
\notready
\end{remark}

% The source section title was: Introduction
Massey used the polar multiplicities of a Lagrangian cycle in the cotangent bundle of a smooth complex manifold to bound the Betti numbers of the stalk of a perverse sheaf at a point \cite[Corollary 5.5]{Massey}. In this paper, we prove an analogous result in characteristic $p$. We use the characteristic cycles for constructible sheaves on varieties of characteristic p defined by \citet[Definition 5.10]{saito1}, building heavily on work of \citet{Beilinson}. Before stating our main theorem, let us define the polar multiplicities. 

\begin{defi}
\source{bounds-stalks-perverse-sheaves-characteristic-p-shende}
\notready We say a closed subset, or algebraic cycle, on a vector bundle is \emph{conical} if it is invariant under the $\mathbb G_m$ action by dilation of vectors. \end{defi}

\begin{defi}
\source{bounds-stalks-perverse-sheaves-characteristic-p-shende}
\notready For a vector bundle $V$ on a variety $X$, let $\mathbb P(V) = \operatorname{Proj} (\operatorname{Sym}^* (V^\vee ))$ be its projectivization, whose dimension $\dim X + \operatorname{rank} V-1$, which is equivalent to the quotient of the affine bundle $V$, minus its zero section, by $\mathbb G_m$. For a conical cycle $C$ on $V$, let $\mathbb P(C)$ the quotient of $C$, minus its intersection with the zero section, by $\mathbb G_m$. \end{defi} 

\begin{defi}
\source{bounds-stalks-perverse-sheaves-characteristic-p-shende}
\notready\label{polar-multiplicity} Let $X$ be a smooth variety of dimension $n$. Let $C$ be a conical cycle on the cotangent bundle $T^* X$ of $X$ of dimension $n$ and let $x$ be a point on $X$.
\source{bounds-stalks-perverse-sheaves-characteristic-p-shende}

For $0 \leq i< \dim X$, let $V$ be a sub-bundle of $T^* X$ defined over a neighborhood of $x$, with rank $i+1$. such that the fiber $V_x$ is a general point of the Grassmanian of $i+1$-dimensional subspaces of $(T^* X)_x$. 

Then we define \emph{the $i$th polar multiplicity of $C$ at $x$},  $\gamma^i_C(x)$, as the multiplicity of the pushforward $\pi_* ( \mathbb P(C) \cap \mathbb P(V))$ at $x$, where $\pi: \mathbb P( T^* X) \to X$ is the projection. 

We define {\emph the $n$th polar multiplicity of $C$ at $x$} to be the multiplicity of the zero-section in $C$. \end{defi}

Here $\pi_*( \mathbb P(C) \cap \mathbb P(V))$ is interpreted as an algebraic cycle, and the multiplicity of an algebraic cycle at a point is the appropriate linear combination of the multiplicities of its irreducible components. We will check that this multiplicity is independent of the choice of $V$ with $V_x$ sufficiently general in Section \ref{equivalences} below.

Our result is as follows:

\begin{theorem}
\source{bounds-stalks-perverse-sheaves-characteristic-p-shende}
\notready\label{first-Massey-intro} Let $X$ be a smooth variety over a perfect field $k$ and let $\ell$ be a prime invertible in $k$. Let $K$ be a perverse sheaf of $\mathbb F_\ell$-modules on $X$. 
\source{bounds-stalks-perverse-sheaves-characteristic-p-shende}

Then $\dim_{\mathbb F_\ell} \mathcal H^{-i}(K)_x$ is at most $i$th polar multiplicity of $CC(K)$ at $x$. \end{theorem}

The analogous statement follows for perverse $\ell$-adic sheaves by noting that their Betti numbers are bounded by the Betti numbers of their mod $\ell$ incarnations.

We have a corollary that describes when these Betti numbers must vanish, which may admit a more direct proof:

\begin{corollary}
\source{bounds-stalks-perverse-sheaves-characteristic-p-shende}
\notready\label{vanishing-intro} Let $X$ be a smooth variety over a perfect field $k$ and let $\ell$ be a prime invertible in $k$. Let $K$ be a perverse sheaf of $\mathbb F_\ell$-modules on $X$. Then $\mathcal H^{-i} (K)_x$ vanishes for \[ -i> \dim SS(K)_x - \dim X \] where $(SS(K))_x$ is the fiber of the singular support of $K$ over $x$. \end{corollary}
\source{bounds-stalks-perverse-sheaves-characteristic-p-shende}

Note that the singular support of a perverse sheaf $K$ is simply the support of its characteristic cycle \cite[Proposiiton 5.14(2)]{saito1}.

Our proof follows to a large extent the strategy of \citep{Massey}. In particular, we successively apply nearby and vanishing cycles to reduce to sheaves on lower-dimensional varieties. However, one key difference is the argument of \citep{Massey} involves passing to an analytic neighborhood and applying vanishing cycles along a sufficiently general analytic function. This general function will in particular be transverse to all the strata of a Whitney stratification of $X$ associated to $K$, except possibly the point $x$ itself. In characteristic $p$, we have access to neither analytic neighborhoods, analytic functions, nor well-behaved Whitney stratifications. Instead, we use a general rational function of degree two (i.e. a pencil of conics). We describe in Lemmas \ref{transversality-check} and \ref{nearby-formula} the transversality conditions the map must satisfy for our argument to work, in terms of the characteristic cycle, and then check, in Lemma \ref{genericity-lemma}, that a general pencil of conics satisfies all these transversality conditions. 

A similar approach can hopefully be used to adapt other arguments involving the characteristic cycle from characteristic $0$ to characteristic $p$.  (Hypersurfaces of larger degree would work equally well, and might come in handy if even more transversality is needed.) 



\subsection{Application to equidistribution in $\Bun_2(\mathbb P^1)$}

In this paper, we prove, as an application of Theorem \ref{first-Massey-intro}:

Let $k$ be a field of characteristic $\neq 2$ and let $C$ be a hyperelliptic curve of genus $g$ over $k$. Define $\Theta_n$ to be the space of degree $n$ effective divisor classes on $C$, viewed as a closed subscheme of the variety $\operatorname{Pic}^n(C)$ parameterizing degree $n$ divisor classes.

\begin{theorem}
\source{bounds-stalks-perverse-sheaves-characteristic-p-shende}
\notready\label{theta-betti-bound} For any $g \in \mathbb N$, $0,\leq a,b \leq g$, and $L \in \operatorname{Pic}^{2g-a-b} (C)$, we have \[ \sum_{i \in \mathbb Z} \dim  H^i (  (\Theta_{g-a} \cap L - \Theta_{g-b} )_{\overline{k}}, \mathbb Q_\ell) \leq 28^g/16 + 4\cdot 8^g + 2\cdot 4^g. \]  \end{theorem}
\source{bounds-stalks-perverse-sheaves-characteristic-p-shende}

This verifies a conjecture of \citet[Conjecture 1.4]{TS}.

\citet[Theorem 4.4]{TS} proved that this conjecture implies a certain equidistribution result, described below:

Let $\operatorname{Bun}_2(\mathbb P^1)$ be the set of isomorphism classes of rank two vector bundles on $\mathbb P^1_{\mathbb F_q}$, up to tensor products with line bundles on $\mathbb P^1_{\mathbb F_q}$. Let $\operatorname{Bun}_2^0(\mathbb P^1)$ be the subset consisting of rank two vector bundles with even degree, and $\operatorname{Bun}_2^1(\mathbb P^1)$ the subset consisting of bundles with odd degree. (Note that these subsets are stable under tensor product with line bundles).  Both $\operatorname{Bun}_2^0(\mathbb P^1)$ and $\operatorname{Bun}_2^1 (\mathbb P^1)$ admit ``uniform" probability measures $\mu_{ \operatorname{Bun}_2^0(\mathbb P^1)} $ and $\mu_{\operatorname{Bun}_2^1 (\mathbb P^1)}$, where the probability of a vector bundle is proportional to the inverse of the order of its automorphism group.  

Let $C$ by a hyperelliptic curve of genus $g$ over $\mathbb F_q$, with a fixed degree two map $\pi: C \to \mathbb P^1$. For $L$ a line bundle on $C$, $\pi_* L$ is a rank two vector bundle on $\mathbb P^1$, and hence defines a point of $\operatorname{Bun}_2(\mathbb P^1)$. This point is preserved by tensoring $L$ with line bundles pulled back from $\mathbb P^1$, so we can think of $\pi_* L$ as a function from $\operatorname{Pic}(C)/ \operatorname{Pic}( \mathbb P^1) $ to $\operatorname{Bun}_2(\mathbb P^1)$. Because  $\operatorname{Pic}(C)/ \operatorname{Pic}( \mathbb P^1) $ is a finite group, it admits a uniform probability measure.

\begin{theorem}
\source{bounds-stalks-perverse-sheaves-characteristic-p-shende}
\notready\label{equidistribution-statement} Let $q> 28^4=  614,656$ be a prime power.
\source{bounds-stalks-perverse-sheaves-characteristic-p-shende}

Fix a sequence of pairs $(C_i,M_i)$ of hyperelliptic curves $C_i$ and line bundles $M_i$ on $C$. Suppose that $\deg M_i \mod 2$ is constant, $g(C_i)$ converges to $\infty$, and the minimum $n$ such that $M_i$ is equivalent in $\operatorname{Pic}(C)/ \operatorname{Pic}( \mathbb P^1) $ to a divisor of degree $n$ converges to $\infty$ with $i$. 

Then as $i$ goes to $\infty$, the pushforward of the uniform probability measure on  $\operatorname{Pic}(C_i)/ \operatorname{Pic}( \mathbb P^1) $ along the map $ L \mapsto (\pi_* L, \pi_* (L \otimes M_i)) $ from $\operatorname{Pic}(C_i)/ \operatorname{Pic}( \mathbb P^1) $  to $\Bun_2(\mathbb P^1)$ converges to 

\[ \frac{1}{2}  \mu_{ \operatorname{Bun}_2^0(\mathbb P^1)} \times \mu_{ \operatorname{Bun}_2^0(\mathbb P^1)} + \frac{1}{2} \mu_{ \operatorname{Bun}_2^1(\mathbb P^1)}\times \mu_{ \operatorname{Bun}_2^1(\mathbb P^1)}\] if $\deg M_i \mod 2=0$ for all i and \[ \frac{1}{2}  \mu_{ \operatorname{Bun}_2^0(\mathbb P^1)} \times \mu_{ \operatorname{Bun}_2^1(\mathbb P^1)} + \frac{1}{2} \mu_{ \operatorname{Bun}_2^0(\mathbb P^1)}\times \mu_{ \operatorname{Bun}_2^1(\mathbb P^1)}\] if $\deg M_i \mod 2 \neq 0$ for all $i$. \end{theorem} 

This follows immediately from Theorem \ref{theta-betti-bound} and \citep[Theorem 4.4]{TS} (which covers in addition the case where the minimum $n$ does does not converge to $\infty$.) 

We now provide some context for these results:

For an imaginary quadratic number field $K$, we can consider the probability measure on the modular curve $X(1)$ that assigns equal measure to the points corresponding to all elliptic curves with complex multiplication by $\mathcal O_K$. Duke's theorem says that, as the discriminant of the fields go to $\infty$, these measures converge to the uniform measure on $X(1)$ \citep{Duke}.

Recalling that, over the complex numbers, there is a natural bijection between the elliptic curves with complex multiplication by $K$ and the class group $Cl(K)$, for each $\alpha$ in $Cl(K)$, let $z_{K,\alpha}$ be point of $X(1)$ of the elliptic curve corresponding to the class group element $\alpha$. For an ideal class $\sigma$, et $\mu_{K,\sigma}$ be the probability measure on $X(1)$ that assigns equal mass to $(z_{K,\alpha}, z_{K, \sigma\alpha})$ for all $\alpha$ in the class group. (One reason this set of points is natural to consider is that it is an orbit under the Galois group $\operatorname{Gal}(\overline{K}|K)$.)

A generalization of Duke's theorem conjectured by \citet[Conjecture 2 on p. 7]{MVUnpublished} is that $\mu_{K,\sigma}$ converges to the uniform measure on $X(1) \times X(1)$ whenever the discriminant of $K$ and the minimal norm of an 
invertible ideal with ideal class $\sigma$ both tend to $\infty$. 

The work of \citet{TS} is a function field analogue of this mixing conjecture. The analogy is constructed by replacing $\mathbb Q$ with $\mathbb F_q(T)$, $X(1)$ with the set $\operatorname{Bun}_2(\mathbb P^1)$, $K$ with the function field of $C$ over $\mathbb F_q$, $Cl(K)$ with $\operatorname{Pic}(C) / \operatorname{Pic}( \mathbb P^1)$, and $z_{K,\alpha}$ with $\pi_* L$. In this setting, Theorem \ref{equidistribution-statement} is exactly the analogue of the conjecture of Michel and Venkatesh (once the trivial but necessary determinant mod $2$ condition is dealt with).

The cohomological conjecture \citep[Conjecture 1.4]{TS} needed to prove this mixing result was proven in characteristic zero by \citet[Theorem 1.5]{TS}, using Massey's bounds for the stalks of perverse sheaves. Thus it was natural to approach the conjecture in characteristic $p$ using Theorem \ref{first-Massey-intro}. Our arguments to prove Theorem \ref{equidistribution-statement} follows closely the proof of \citep[Theorem 1.5]{TS}. One modification needed is that, in characteristic zero, one knows that all irreducible components of the characteristic cycle are Lagrangian varieties, hence are the conormal bundle to their supports, and one can thus calculate the characteristic cycle by studying only these supports. In characteristic $p$, these irreducible components are not necessarily Lagrangian, and so it is necessary to perform calculations in the cotangent bundle, not on the base variety. In addition, the argument uses a new idea provided by Tsimerman in the appendix to bound a crucial multiplicity. 




Since the writing of \citep{TS}, the equidistribution conjecture on $X(1) \times X(1)$ was verified by \citet[Theorem 1.3]{Khayutin}, using ergodic theory methods. In addition, \citet{Khayutin} proved this statement over modular curves of higher level (while Theorem \ref{equidistribution-statement} requires level $1$.) However, this required two additional assumptions: that the fields $K$ are always split at two fixed primes $p_1, p_2$, and that their Dedekind zeta functions have no Landu-Siegel zero.  

In comparing these results, one should note that (unlike some results over $\mathbb Q$) it is not yet clear if the argument of \citet{Khayutin} can be made to work over function fields, as there are more measures to rule out. See \citep*[Theorem 1.2 and \S1.3]{ELM} for a measure classification result and a discussion of the difficulties arising from measures invariant under subgroups defined over subfields, of which $\mathbb F_q(T)$ has infinitely many. Such a transfer would allow one to remove the level 1 assumption from Theorem \ref{equidistribution-statement}, at the cost of introducing the split primes assumption. Going from the function field to the number field case, on the other hand, is as hard as usual.



\subsection{Acknowledgments}

The author was supported by Dr. Max R\"{o}ssler, the Walter Haefner Foundation and the ETH Zurich Foundation, and, later, served as a Clay Research Fellow, while working on this research. I would like to thank Takeshi Saito, Vivek Shende, and Jacob Tsimerman for helpful discussions about their works, Philippe Michel and Manfred Einsiedler for helpful comments about the general equidistribution problem, and the anonymous referee for helpful comments.
